\section{From the UOR Atlas to Its Carried Modular Category}
\label{sec:bridge}

This section connects the object the reader knows from the UOR program---the Atlas---to the modular tensor category realized on top of it, so that the constructions of the later sections are read as properties of a single, already-defined structure rather than as free-standing definitions.

\subsection{The Atlas as a parametric, content-addressed object}
The UOR Atlas is fixed by the parametric tuple $(q, T, O) = (4, 3, 8)$: a \emph{scope} $q$, a \emph{modality} $T$, and a \emph{context} $O$. Its label space is the set of $q\cdot T\cdot O = 96$ classes, and its identity mechanism is \emph{content addressing}: every state, generator, and evaluated braid is named by a Universal Object Reference $\kappa$---a cryptographic address of its canonical bytes on the \textit{holospaces} substrate~\cite{holospaces}. Two computations carry the same $\kappa$ if and only if they have the same content, so isotopic braid words (distinct paths composing to one topological operator) name the same $\kappa$; this is the identity relation the framework computes over, and the mechanism behind the measured deduplication of Section~6. All Atlas constants used here---the class count and stride, the generator orders, the balanced-operator spectrum, the $E_8$ seed, the modular and Coxeter data---are drawn from the machine-checked Lean~4 formalization F1 as an external oracle, pinned by digest and upstream tag; the dependency runs one way (the UTQC validates against F1; F1 does not depend on the UTQC).

\subsection{The $24$-class carrier}
The $96$ labels are a $\mathbb{Z}_q$-graded extension of a $24$-dimensional carrier: quotienting by the scope grading $q = 4$ leaves the $24 = T\cdot O$ classes that index the carrier $V_T \otimes V_O$ (with $T = 3$, $O = 8$). These $24$ classes are indexed by a set of the cardinality of an $S_4$-orbit, $|S_4| = 24$, and are acted on by the three reflection generators $(\sigma, \tau, \mu)$ of orders $(q, O, 2) = (4, 8, 2)$; they are an index set, not the group $S_4$ (whose elements have no order~$8$, unlike $\tau$). The balanced spectral operator on this carrier has the sourced spectrum $\{10, 7, 2, -1\}$ with multiplicities $\{1, 2, 7, 14\}$ and signature $(10, 14)$, which the implementation derives parametrically and cross-checks against F1.

\subsection{Why this structure carries a pointed MTC}
The $24$-class carrier is exactly the object set of an abelian group: $A = \mathbb{Z}_{T} \times \mathbb{Z}_2^{3} = \mathbb{Z}_3 \times \mathbb{Z}_2^{3}$, the product of a $\mathbb{Z}_3$ anyon with three semions ($|A| = 24$). Group-law fusion on $A$---the non-negative associative quotient of the signed octonion composition that the substrate realizes---makes the carrier a pointed fusion category, and the quadratic form $q(m, c) = \zeta_3^{m^2}\, i^{|c|}$ equips it with the Eilenberg-MacLane abelian $3$-cocycle as associator and braiding. A pointed fusion category with a nondegenerate quadratic form is a pointed \emph{modular} tensor category, $\mathcal{C}(A, q)$; Section~\ref{sec:atlas_analysis} verifies its modular axioms phase-exactly, and Section~3 decides the density of the coupled generators on its carriers. The remainder of the paper is the analysis of this one category $\mathcal{C}(A, q)$ carried by the Atlas, addressed throughout by $\kappa$.
